\section{Введение}
\label{sec:Chapter0} \index{Chapter0}

%В этой части надо описать предметную область, задачу из которой вы будете
% решать, объяснить её актуальность (почему надо что-то делать сейчас?). Здесь же
% стоит ввести определения понятий, которые вам понадобятся в постановке задачи.

Современные вычислительные системы всё чаще строятся на основе архитектуры с неоднородным доступом к памяти (NUMA, Non-Uniform Memory Access), в которой время доступа к памяти зависит от расположения вычислительного узла (процессора) относительно физической памяти. В NUMA-системах память логически разделена между нодами, и доступ к локальной памяти осуществляется быстрее, чем к удалённой (см рис.~\ref{fig:numa_basic}). Такая организация позволяет масштабировать производительность на многоядерных и многопроцессорных системах, однако одновременно приводит к возникновению так называемого NUMA-эффекта~--- деградации производительности при частом обращении потоков к удалённой памяти.

\begin{figure}[H]
    \centering
    \includegraphics[width=1\textwidth]{NUMA_access_latency.pdf}
    \caption{Задержка доступа в зависимости от локальности}
    \label{fig:numa_basic}
\end{figure}

Операционные системы, в частности, ядро Linux~\cite{linux}, предоставляют механизмы управления размещением памяти и потоков, включая политики NUMA и средства привязки. Идеальная ситуация, которую пытаются достичь такие механизмы,~--- когда все потоки находятся на одной NUMA-ноде со своей памятью, и, как следствие, все доступы локальны (см рис.~\ref{fig:numa_ideal}).

Тем не менее, в реальных пользовательских приложениях часто наблюдается сложное и динамическое соответствие между потоками и используемой памятью. Это особенно характерно для высоконагруженных серверных приложений, систем обработки данных и многопоточных вычислений, где доступ к памяти может быть распределён неравномерно и изменяться во времени. В таких условиях традиционные подходы управления размещением памяти оказываются недостаточно эффективными, и проблема NUMA-эффекта остаётся актуальной.

\begin{figure}[H]
    \centering
    \includegraphics[width=1\textwidth]{NUMA_local_access.pdf}
    \caption{Идеальная ситуация~--- все доступы локальны}
    \label{fig:numa_ideal}
\end{figure}

\hypertarget{par:NUMA_repl_intro}{}
В данной работе рассматривается новый подход, направленный на снижение влияния NUMA-эффекта за счёт репликации таблиц трансляций и частей адресного пространства процесса. Ключевая идея заключается в создании локальных копий выбранных областей памяти для каждой NUMA-ноды. Это позволит потокам приложения, выполняющимся на конкретном узле, преимущественно обращаться к локальным копиям, тем самым снижая долю удалённых обращений к памяти (см рис.~\ref{fig:numa_replication}).

\begin{figure}[H]
    \centering
    \includegraphics[width=1\textwidth]{NUMA_replication_idea.pdf}
    \caption{Идея механизма репликации}
    \label{fig:numa_replication}
\end{figure}
